% Part: computability
% Chapter: recursive-functions
% Section: composition

\documentclass[../../../include/open-logic-section]{subfiles}

\begin{document}

\olfileid[fr]{cmp}{rec}{com}
\olsection{Composition}

Si $f$ et $g$ sont deux fonctions des entiers naturels à un
argument, on peut les composer : $h(x)=f(g(x))$. La nouvelle
fonction~$h$ est alors définie par \emph{composition} à partir
de $f$ et de~$g$. Généralisons ce procédé aux fonctions de
plusieurs arguments.

Supposons que $f$ soit une fonction à $k$ arguments et que
$g_0,\dots,g_{k-1}$ soient $k$ fonctions, chacune à $n$
arguments. On peut définir une fonction~$h$ à $n$ arguments par :
\[
h(x_0, \dots, x_{n-1}) =
f(g_0(x_0, \dots, x_{n-1}), \dots, g_{k-1}(x_0, \dots, x_{n-1}))
\]
Si $f$ et toutes les fonctions $g_i$ sont calculables, $h$ l'est
aussi. Pour calculer $h(x_0,\dots,x_{n-1})$, on calcule d'abord,
pour chaque $i=0,\dots,k-1$, la valeur
$y_i=g_i(x_0,\dots,x_{n-1})$. On donne ensuite ces valeurs à~$f$
pour obtenir $h(x_0,\dots,x_{n-1})=f(y_0,\dots,y_{k-1})$.\footnote{Correction éditoriale : les arguments de $h$ sont indexés jusqu'à $n-1$, et non $k-1$ comme dans cette phrase de la source. $n$ est l'arité du résultat ; $k$ est celle de $f$.}

Cette description peut sembler trop restrictive pour rendre compte
des façons dont on utilise des fonctions connues pour en calculer
de nouvelles. Il arrive, d'abord, que l'on n'utilise pas tous les
arguments d'une fonction : c'est le cas de $g(x,y,z)=\Succ(z)$,
employée dans la définition par récurrence primitive de~$\Add$.
Supposons que l'on dispose des fonctions suivantes :
\[
\Proj{n}{i}(x_0, \dots, x_{n-1}) = x_i
\]
On les appelle \emph{fonctions de projection} : $\Proj{n}{i}$
est une fonction à $n$ arguments, avec $n\geq1$ et $0\leq i<n$.\footnote{Correction éditoriale : nous écrivons $\Proj{n}{i}$ dans la phrase où la source emploie $\Proj{k}{i}$. Les indices commencent à $0$ dans OpenLogic, contrairement à ceux du témoin français.} On peut alors définir $g$ par :
\[
g(x, y, z) = \Succ(\Proj{3}{2}(x, y, z)).
\]
Ici, le rôle de $f$ est tenu par la fonction à un argument
$\Succ$, de sorte que $k=1$. La fonction à trois arguments
$\Proj{3}{2}$ joue celui de~$g_0$. Le résultat est une fonction
à trois arguments qui renvoie le successeur du troisième.

Les projections permettent également de définir de nouvelles
fonctions en permutant ou en identifiant des arguments. Ainsi,
on peut définir la fonction $h(x)=\Add(x,x)$ par :
\[
h(x_0) = \Add(\Proj{1}{0}(x_0),\Proj{1}{0}(x_0)).
\]
Ici $k=2$ et $n=1$ : $\Add$ joue le rôle de $f(y_0,y_1)$,
tandis que les rôles de $g_0(x_0)$ et de $g_1(x_0)$ sont tous deux
tenus par $\Proj{1}{0}(x_0)$, projection à un argument, autrement
dit la fonction identité.

Si l'on dispose d'une fonction $f(y_0,y_1)$, on peut définir
$h(x_0,x_1)=f(x_1,x_0)$ par :
\[
h(x_0, x_1) = f(\Proj{2}{1}(x_0, x_1),\Proj{2}{0}(x_0, x_1)).
\]
Ici $k=n=2$, et $\Proj{2}{1}$ et $\Proj{2}{0}$ jouent
respectivement les rôles de $g_0$ et de~$g_1$.

On pourrait aussi craindre que l'obligation d'attribuer la même
arité~$n$ à $g_0,\dots,g_{k-1}$ ne limite ce procédé. Rappelons
que l'\emph{arité} d'une fonction est son nombre d'arguments :
une fonction à $n$ arguments est d'arité~$n$. Les projections
apportent pourtant la souplesse voulue. Supposons, par exemple,
que $f$ et $g$ aient chacune trois arguments et que $h$ soit la
fonction à deux arguments définie par :
\[
h(x,y) = f(x,g(x,x,y),y).
\]
À l'aide des projections, on peut réécrire cette définition :
\[
h(x,y) = f(\Proj{2}{0}(x,y), g(\Proj{2}{0}(x,y), \Proj{2}{0}(x,y),
\Proj{2}{1}(x,y)), \Proj{2}{1}(x,y)).
\]
Ainsi, $h$ est obtenue par composition de $f$ avec
$\Proj{2}{0}$, $l$ et $\Proj{2}{1}$, où
\[
l(x,y) = g(\Proj{2}{0}(x,y),\Proj{2}{0}(x,y),\Proj{2}{1}(x,y)),
\]
c'est-à-dire que $l$ est elle-même obtenue par composition de $g$
avec $\Proj{2}{0}$, $\Proj{2}{0}$ et~$\Proj{2}{1}$.

\end{document}
